% Appendix A, typeset from info/simulation_workflow.md.

\chapter{Simulation with GHDL}
\label{app:sim}

This is the book's permanent reference for simulation. Every chapter that verifies a module against a
testbench points here instead of repeating the steps. \Secref{c2:sec:testbenches} is the gentler
introduction, written for the first time you run a testbench at all; this appendix is the one you come
back to, and it covers the group project's build in particular.

% ------------------------------------------------------------------------------------------
\section{Why simulate}
\label{app:sim:why}

\Chapref{c1:ch} to \ref{c8:ch} verify every design by hand, in CircuitVerse or directly on the
DE0-CV. That works for small circuits with a handful of inputs a human can toggle.

It stops working as soon as a design's correctness depends on exact timing relationships across
several cycles between several modules, which is the group project's situation throughout.

\begin{itemize}
  \item A \textbf{testbench} is a small VHDL program, handed out, that drives a design's inputs and
    checks its outputs automatically.
  \item \textbf{GHDL} is the tool that compiles and runs it, entirely in software, with no FPGA
    involved.
  \item You are never asked to write the testbench that \emph{verifies} a module you build. Every one
    of those is handed out; you run it and read what it reports.
  \item A few exercises invite you to write a short \emph{ad hoc} rig, twenty lines feeding your own
    bit sequence into a module and \code{report}ing what comes out. That is exploratory, and always
    optional in form: nothing you are assessed on depends on a testbench you wrote yourself.
\end{itemize}

% ------------------------------------------------------------------------------------------
\section{Installing GHDL}
\label{app:sim:install}

On Ubuntu or WSL:

\begin{shell}
sudo apt-get update && sudo apt-get install -y ghdl
\end{shell}

Confirm that it worked:

\begin{shell}
ghdl --version
\end{shell}

GHDL 3.x or later. The course's CI runs the same GHDL commands as below, so what passes locally passes
in CI too.

% ------------------------------------------------------------------------------------------
\section{Where the files live}
\label{app:sim:layout}

The group project's VHDL lives in \textbf{two flat directories at the repository root}: the modules
you write, alongside the testbenches handed out (\file{*_tb.vhd}) that check them. There are no
per-session directories and nothing to copy between them; \file{can_def.vhd} in particular is handed
out, exists \emph{once}, and every controller module reads that one file.

The split between the two directories is the protocol boundary: everything in \file{controller/}
knows something about CAN and nothing about SPI, and everything in \file{bridge/} the other way
round.

\begin{textdrawing}
controller/   can_controller.vhd     ch 11, yours (ports only at first, grows through ch 17)
              meta_prev.vhd          ch 12, yours
              bit_timer.vhd          ch 12, yours
              crc15.vhd              ch 13, yours
              tx_shift_reg.vhd       ch 14, yours
              rx_shift_reg.vhd       ch 15, yours
              + can_def.vhd and the six *_tb.vhd, handed out

bridge/       register_bank.vhd      ch 18, yours
              spi_reg_bridge.vhd     ch 19, yours
              can_spi_node.vhd       ch 19, yours
              + spi_slave.vhd, spi_def.vhd, register_bank_tb.vhd,
                spi_reg_bridge_tb.vhd, handed out
\end{textdrawing}

The two directories initially contain only the files handed out: nine testbenches, plus
\file{can_def.vhd}, \file{spi_slave.vhd} and \file{spi_def.vhd}. Every chapter from \chapref{c11:ch}
adds a module, with
the exception that \file{can_controller.vhd} is written early (\chapref{c11:ch}) as an entity with an
empty architecture and then \emph{grows}: \chapref{c16:ch} fills in its transmit path and
\chapref{c17:ch} its receive path, without any new file appearing for either. That is why
\code{can_controller_tb}, the only testbench exercising the whole controller, is gated on more than
the files existing; \secref{app:sim:make} returns to that.

% ------------------------------------------------------------------------------------------
\section{The three GHDL steps}
\label{app:sim:steps}

Every VHDL design (and testbench) goes through the same three steps, in this order:

\begin{enumerate}
  \item \textbf{Analyse} (\code{ghdl -a}): parses a \file{.vhd} file, checks it for syntax and type
    errors, and registers its design units (entities, architectures, packages) in a work library.
    Dependencies have to be analysed before the units that use them: \file{can_def.vhd} before every
    module that says \code{use work.can_def.all;}, and a module before its own testbench.
  \item \textbf{Elaborate} (\code{ghdl -e}): resolves a particular top-level entity's complete
    instantiation hierarchy into something executable.
  \item \textbf{Run} (\code{ghdl -r}): actually simulates it.
\end{enumerate}

For a module with a testbench handed out, for example \chapref{c12:ch}'s \code{meta_prev}, which
reads no package at all and therefore needs only its own two files:

\begin{shell}
cd controller
ghdl -a --std=93 meta_prev.vhd meta_prev_tb.vhd
ghdl -e --std=93 meta_prev_tb
ghdl -r --std=93 meta_prev_tb --assert-level=error
\end{shell}

Every other controller module reads \code{can_def}, so the package comes first:

\begin{shell}
cd controller
ghdl -a --std=93 can_def.vhd bit_timer.vhd bit_timer_tb.vhd
ghdl -e --std=93 bit_timer_tb
ghdl -r --std=93 bit_timer_tb --assert-level=error
\end{shell}

\file{bridge/} works the same way, with one wrinkle: \code{register_bank} reads \code{can_def}, which
is a directory away.

\begin{shell}
cd bridge
ghdl -a --std=93 ../controller/can_def.vhd register_bank.vhd register_bank_tb.vhd
ghdl -e --std=93 register_bank_tb
ghdl -r --std=93 register_bank_tb --assert-level=error
\end{shell}

\begin{itemize}
  \item \code{--std=93} selects VHDL-93 throughout the course. Every testbench handed out terminates
    itself cleanly, by setting a \code{done} signal the clock generation process checks and then
    calling \code{wait;}, rather than relying on any ``stop the simulator'' procedure from a later
    standard.
  \item \code{--assert-level=error} tells GHDL to stop the simulation, with a non-zero exit code, at
    every failed check reported at level \code{error} or higher. \textbf{Always pass it when you run a
    project testbench.} The next section says why it is not optional here.
\end{itemize}

\subsection{Why \code{--assert-level=error} matters, and where it does not}
\label{app:sim:assert}

GHDL stops on a failed assertion only when the assertion's severity reaches the assert level, and
GHDL's default level is \code{failure}. The book's two halves sit on opposite sides of that boundary,
which is worth knowing since it explains why the flag looks optional in part~\ref{part:vhdl} and is
not in part~\ref{part:can}:

\begin{booktable}{lcLL}
  & \thead{Level} & \thead{Without the flag} & \thead{With the flag} \\
  \midrule
  The chapters' testbenches & \code{failure} & stops, exits 1 & stops, exits 1 \\
  The project's testbenches & \code{error} & \textbf{prints, carries on, exits 0} & stops, exits 1 \\
\end{booktable}

For a project testbench, then, omitting the flag turns a real failure into a message scrolling past
above a zeroed exit code, which is exactly what a ``silent pass'' looks like. Pass it every time and
the difference stops mattering.

\code{--stop-time=10ms} is worth adding for the same kind of reason. A testbench waiting for an event
your module never produces does not \emph{fail}; it hangs. The stop time turns that into a run that
ends. It is not in the commands above because the project's testbenches all bound their own waits, but
it costs nothing and \secref{c2:sec:testbenches} makes it a habit.

Every testbench reports only pass or fail (and, on failure, the reason), nothing else. A passing run
ends with a line like this:

\begin{console}
bit_timer_tb.vhd:148:9:@7312ns:(report note): bit_timer: all checks passed!
\end{console}

A failing run stops early with an \code{(assertion error)} message naming exactly which check failed
and why. Read that message; it is written to say specifically what went wrong.

% ------------------------------------------------------------------------------------------
\section{Building everything at once}
\label{app:sim:make}

\code{make build-project}, from the repository root, does all of this automatically for every
testbench whose modules you have written. It is useful for confirming that your whole tree still
passes, but it is the understanding of the three individual steps above that lets you debug a single
module quickly during a session.

Since the project is built up one module at a time, \code{make build-project} \textbf{skips} every
testbench whose modules do not exist yet, instead of failing on them, and says so:

\begin{console}
==> meta_prev_tb
controller/meta_prev_tb.vhd:114:9:@311ns:(report note): meta_prev: all checks passed!
==> bit_timer_tb (skipped: no bit_timer.vhd yet)
\end{console}

That is what \chapref{c12:ch} looks like halfway through: \file{meta_prev.vhd} written and passing,
\file{bit_timer.vhd} still to do. \file{can_def.vhd} is never named as missing, since it is handed
out.

Skipped is not failed, deliberately. The project runs over ten sessions and most of the tree is
missing during most of them; a build glowing red until the last module landed would say nothing useful
during any of the previous nine. The number of testbenches actually running therefore grows as you go:
none after \chapref{c10:ch} or \ref{c11:ch}, two after \chapref{c12:ch}, then one more in each of
\chapref{c13:ch}, \ref{c14:ch} and \ref{c15:ch}, one more in \chapref{c17:ch}, one in
\chapref{c18:ch}, and the last in \chapref{c19:ch}.

\begin{aside}
\code{can_controller_tb} is \textbf{gated on more than the file existing}, and that is worth knowing
before it puzzles you. \file{can_controller.vhd} is written in \chapref{c11:ch} as an entity with an
empty architecture and grows through \chapref{c16:ch} and \ref{c17:ch}, so from \chapref{c15:ch}
onwards every file the testbench needs exists while the controller still cannot receive a frame.
Running it then would fail, twice over: once in \chapref{c15:ch}, and again in \chapref{c16:ch}, which
can transmit but not yet receive. Since the testbench works by having one node receive another's
frame, what it actually requires is a \textbf{receive path}, and that comes in \chapref{c17:ch}. The
build script therefore looks inside your \file{can_controller.vhd} for \chapref{c17:ch}'s
receive-side CRC gating and reports:

\begin{console}
==> can_controller_tb (skipped: can_controller.vhd has no receive path yet;
    set CI_BUILD_ALL=1 to run it anyway)
\end{console}

That check is a heuristic rather than a real dependency, and it has a failure mode worth naming: if
your \chapref{c17:ch} is correct but spells those expressions differently from the chapter, you get
this skip message when the testbench should have run. Nothing is broken; run
\code{CI_BUILD_ALL=1 make build-project} to force every testbench regardless. The check cannot fail the
other way round, since the marker cannot appear before a receive path does.
\end{aside}

% ------------------------------------------------------------------------------------------
\section{Cleaning up}
\label{app:sim:clean}

GHDL's analysis step creates \file{work-obj93.cf} and similar files in the directory you ran them
from. The build script keeps a \file{work/} directory of its own instead, so those files do not
litter the source tree. If you run the three steps by hand in \file{controller/} or \file{bridge/},
clean up with:

\begin{shell}
make clean
\end{shell}

which removes the generated \file{work/} directories, any stray \file{work-obj*.cf}, and the waveform
and executable files a hand run leaves behind.
